% Proved after four actual-file proof audits; exact candidate and source-read scopes retained.
\subsection{Field-semilinear counts and unramified field-permutation crowns}

We retain the complete one-smooth-exceptional-curve setting of
Theorem~\ref{D:etale-resolution-hereditary-linear-count}. In this
subsection $C^2=-\delta<0$, $g=g(C)$ and $\Gamma=\pi_1(U)$.
All representations are continuous. For $q=p^f$, write
$\mathcal G_q=\operatorname{Gal}(\mathbb F_q/\mathbb F_p)$.
It acts entrywise on $\operatorname{GL}_r(\mathbb F_q)$.

\begin{lemma}[Proved: twisted Frobenius models of a fixed field character]
\label{D:twisted-deck-Frobenius-semilinear-model}
Let $W$ be a connected locally noetherian scheme over $\mathbb F_q$
with a geometric base point, and fix a
surjective continuous character
$\chi:\pi_1(W)\to\langle\sigma\rangle\subset\mathcal G_q$,
where $\sigma(a)=a^{q_0}$, $q_0=p^e$, $f=ed$ and
$|\langle\sigma\rangle|=d$. Let $v:V\to W$ be its connected
cyclic finite etale cover. Choose the generator $\tau$ of its
deck transformations with the pullback convention for which a
descent map $\tau^*\mathcal L\to\mathcal L$ is
$\sigma$-semilinear. This is the inverse of the corresponding
monodromy permutation convention.

The groupoid of rank-$r$ representations
\[
 \rho:\pi_1(W)\longrightarrow
      \operatorname{GL}_r(\mathbb F_q)\rtimes\langle\sigma\rangle
 \quad\hbox{with field projection }\chi,
                                                        \tag{D.SC.1}
\]
with conjugacies by $\operatorname{GL}_r(\mathbb F_q)$, is
equivalent to the groupoid of vector bundles $E$ on $V$ with
an isomorphism
\[
             \Phi:\tau^*F_{q_0}^*E\xrightarrow{\sim}E.  \tag{D.SC.2}
\]
Here $F_{q_0}$ is absolute Frobenius; a morphism is a bundle
map commuting with (D.SC.2). In the relative version, finite
local systems replace the chosen fundamental-group notation.
That construction works after arbitrary $\mathbb F_q$-coefficient-base
extension. In the surface applications all schemes and coefficient
algebras are over the fixed algebraically closed field $k$.
\end{lemma}

\begin{proof}
Put $T=\tau^*F_{q_0}^*$. The two pullbacks commute, and
$\tau^d=1$, so there is a canonical identification
$T^d=F_q^*$. Define, rather than independently prescribe,
\[
 \Psi=\Phi\circ T\Phi\circ\cdots\circ T^{d-1}\Phi:
                              F_q^*E\xrightarrow{\sim}E.\tag{D.SC.3}
\]
The ordinary unit-bundle/Lang-frame equivalence used in
Proposition~\ref{D:finite-jet-unit-Isom} gives a locally
constant sheaf $\mathcal L$ of rank-$r$ $\mathbb F_q$-vector
spaces on $V$. Concretely its sections are the vectors fixed
by $\Psi$; its frame functor is finite etale and locally
nonempty. This construction and its inverse, extension of
scalars from $\mathbb F_q$ to $\mathcal O_V$, work over any
coefficient algebra over $\mathbb F_q$. No perfectness or
reducedness of that algebra is required: the equations for
Lang frames have invertible linear derivative.

The identity obtained by rotating the $d$ factors in
(D.SC.3) shows that $\Phi$ carries fixed frames to fixed
frames. It induces a $\sigma$-semilinear map
$D:\tau^*\mathcal L\to\mathcal L$. Its $d$-fold deck
composite is the action of $\Psi$ on its fixed sections,
and is therefore the identity. Thus $D$ is genuine cyclic
semilinear descent, with the prescribed field character.
Effective finite etale descent gives (D.SC.1).

Conversely, a representation (D.SC.1) restricts on
$\ker\chi$ to an $\mathbb F_q$-linear local system
$\mathcal L$ on $V$, with this semilinear deck descent.
Set $E=\mathcal O_V\otimes_{\mathbb F_q}\mathcal L$.
The $\sigma$-semilinearity is exactly the scalar twist
needed to extend $D$ to the linear map (D.SC.2) from
$\tau^*F_{q_0}^*E$. Its $d$-fold composite is the ordinary
$F_q$-unit map of $\mathcal L$. These constructions are
inverse on objects and on the stated morphisms, since
frames recover all the descent maps.

This is a categorical deck-cocycle assertion. It does not
assert that an arbitrarily chosen element of $\pi_1(W)$
lifting $\sigma$ has order $d$, or that its semilinear
matrix has $d$th power one. The $d$th power of that element
lies in $\pi_1(V)$; its path transport is already present
in the descended local system. Replacing it by a trivial
pointed group relation would lose representations.
\end{proof}

\begin{theorem}[Proved: an index-linear count with unramified field action]
\label{D:etale-resolution-semilinear-count}
There is a constant $A_R'$ with the following property.
Let $Y_H\to Y$ be a connected finite etale cover of degree
$t$, set $U_H=Y_H\times_Y U$ and $H=\pi_1(U_H)$, and
let $\chi:H\to\mathcal G_q$ be any character factoring
through $\pi_1(Y_H)$. For every $r\geq1$, the number of
representations
\[
 H\longrightarrow\operatorname{GL}_r(\mathbb F_q)
                         \rtimes\mathcal G_q
 \quad\hbox{with field projection }\chi,
                                                        \tag{D.SC.4}
\]
up to conjugation by $\operatorname{GL}_r(\mathbb F_q)$,
is at most $q^{A_R'tr^2}$. The constant is independent of
$q,t,r,\chi$ and the order of its field image. The assumption
that $\chi$ extends over $Y_H$ is part of the statement.
\end{theorem}

\begin{proof}
Write $|\operatorname{im}\chi|=d$, $e=f/d$ and $q_0=p^e$.
The character defines a connected cyclic degree-$d$ finite
etale cover $Y_V\to Y_H$, with deck generator $\tau$,
and $V=Y_V\times_Y U$. Its composite degree over $Y$
is $td$. As in the proof of
Theorem~\ref{D:etale-resolution-hereditary-linear-count},
the Stein ring is a complete normal local surface ring,
and $Y_V$ is a projective regular resolution with sole
smooth exceptional curve $C_V$. We have
\[
       g(C_V)-1=td(g-1),\qquad C_V^2=-td\delta.          \tag{D.SC.5}
\]
The deck map preserves $C_V$ and its normal ideal.

By Lemma~\ref{D:twisted-deck-Frobenius-semilinear-model},
we count (D.SC.2). Factor it into $e$ one-step
$p$-Frobenius arrows, placing $\tau^*$ on the last arrow.
For instance the intermediate punctured bundles may be
$F_p^{i*}E$ for $0\leq i<e$; allowing all isomorphic
one-step presentations gives the same groupoid. This uses
$e$ stages, not $f$ stages on the degree-$td$ cover.

\emph{The determinant is counted before passing to $V$.}
Taking the determinant in (D.SC.4) gives a multiplicative
cocycle in the finite $H$-module
\[
           A_\chi=\mathbb F_q^*,\qquad
                    h\cdot a=\chi(h)(a).
\]
Its isomorphism class lies in $H^1(H,A_\chi)$. The
twisted-coefficient bound of Lemma~\ref{D:unramified-semilinear-determinant-budget}
for this regular one-divisor pair gives
\[
 |H^1(H,A_\chi)|
       \leq(q-1)^{2g(C_H)+1}\leq q^{(2g+3)t}.          \tag{D.SC.6}
\]
For clarity, the ingredients are that $A_\chi$ extends
to a finite local system on $Y_H$, the low-degree
localization sequence has cokernel in
$H^0(C_H,A_\chi(-1))$, and proper henselian base change
reduces $H^1(Y_H,A_\chi)$ to the smooth proper curve
$C_H$. Its fundamental group has at most $2g(C_H)$
generators. This reasoning allows $p\mid d$ because
the coefficient group $A_\chi$ itself has order prime
to $p$; the field-action monodromy is unramified on $Y_H$.
No list of arbitrary ordinary characters on the larger
cover $V$ is used.

Fix one determinant class in (D.SC.6). Its restriction to
$V$ fixes the ordinary punctured determinant line of $E$.
At stage $i$ it fixes its $p^i$th pullback. All extensions
of one such line to $Y_V$ differ by
$\mathcal O_{Y_V}(bC_V)$, $b\in\mathbb Z$. The minimal
residual lattice window has degree width $rtd\delta$,
whereas a change by $C_V$ changes degree by $td\delta$.
There are therefore at most $r+1$ possible determinant
extension lines at each stage, or $(r+1)^e$ lists.

\emph{Uniform actual families with $e$ stages.}
The fixed-line construction in the proof of
Theorem~\ref{D:etale-resolution-hereditary-linear-count}
applies with the cover degree there replaced by $td$.
Its ordinary pole bounds and normal precision $a$ are
independent of that degree: both the genus term and the
positive normal degree in the required slope inequalities
scale by $td$ in (D.SC.5). The same inequalities hold at
the last arrow with $\tau^*$, because it preserves slopes,
degrees, the normal ideal and the conormal twists.
Closing the slope recurrence around this permuted cycle
uses that preservation, not the assertion $\tau=1$.

After the same fixed ordinary twist, each residual bundle
is a quotient of $\mathcal O_{C_V}^{r+1}$ with the fixed
determinant kernel line. Its generators lift through the
chosen normal precision. Applying the actual filtered
Cech construction and the finite quotient-Hom complexes
gives at most $A\,td\,r^2$ coordinates per one-step stage.
Only $a$ ordinary normal layers occur. Reference line
cochains and their linear cohomology splittings are fixed
coefficients for each discrete determinant list.

We can use the pullbacks of two fixed affine complements
on $C$ for these Cech computations; they are affine and
are preserved by $\tau$. Thus the last pullback is an
actual linear operation on the fixed complexes. It adds
neither variables nor factors of variable degree. The
normal perturbation words have length at most $a$,
internal width two, fixed-length compositions and only
one $p$th power per arrow. The complete polynomial
presentation consequently has
\[
     N\leq A\,e\,td\,r^2=A ftr^2,
       \qquad \deg(\text{each equation})\leq D_0,       \tag{D.SC.7}
\]
with $A,D_0$ fixed for $R,p$. Its open conditions are the
same actual subbundle, slope and inverse conditions as
in the fixed-line theorem. This is an all-coefficient
presentation of objects and arrows, not a tangent bound.

The proof of Proposition~\ref{D:finite-jet-unit-Isom}
also applies to this final deck pullback. In the normal
contraction estimate, $\tau^*$ preserves order, so the
one-step inequality remains $j\mapsto pj-2B$. The map
and inverse jets, their specified lifting, and their
algebraization use the same fixed ordinary precision.
After arbitrary Artinian coefficient extension, the
Lang frames of (D.SC.3), together with their finite
semilinear deck descent, are finite etale. Nilpotent
invariance therefore gives exactly the same uniqueness
and lifting of isomorphisms. The represented class
fibres are open and closed. Notice that (D.SC.3) is
used for effectivity, not expanded as an equation of
degree growing with $d$ in (D.SC.7).

By Lemma~\ref{D:weighted-affine-component-bound}, the
number of unit classes in each such open family is at
most $(N+1)D_0^N\leq\exp(A_1ftr^2)$. Multiply by
$(r+1)^e$ and (D.SC.6). Since $e\leq f$, $t\geq1$
and $\log(r+1)=O(r^2)$, their product is
$q^{A_R'tr^2}$ after enlarging a constant. Lemma
\ref{D:twisted-deck-Frobenius-semilinear-model} identifies
the counted isomorphisms with the conjugacies stated in
(D.SC.4). This proves the claim, also when $d=1$.
\end{proof}

\begin{theorem}[Proved: bounded generators for unramified field-permutation crowns]
\label{D:unramified-field-permutation-crowns}
Fix $\kappa>0$. Let $S$ be a nonabelian finite simple
group, embedded faithfully as a subgroup of
$\operatorname{GL}_r(\mathbb F_q)$ stable under entrywise
$\mathcal G_q$, and assume
\[
                    r^2\log q\leq\kappa\log|S|.       \tag{D.SC.8}
\]
Let $P\leq\mathcal G_q^t\rtimes\operatorname{Sym}(t)$
act on $A=S^t$ by its field automorphisms and permutations.
Assume the permutation action is transitive and the
induced map $P\to\operatorname{Out}(A)$ is injective.
Set $L=A\rtimes P$, with socle $A$. If a continuous
surjection
\[
              \Gamma\twoheadrightarrow L_k=A^k\rtimes P
                                                        \tag{D.SC.9}
\]
has its quotient $\beta:\Gamma\twoheadrightarrow P$
factoring through $\pi_1(Y)$, then, for
$c=\max\{1,\lceil\kappa(A_R'+1)\rceil\}$,
\[
             k\leq|A|^c,
        \qquad d(L_k)\leq\max\{2,2g,c+2\}.            \tag{D.SC.10}
\]
The same bound works for all displayed $S,q,t,k,P,\beta$.
This includes field automorphisms of unbounded order.
It excludes quotients $\beta$ ramified over the exceptional
curve and outer actions not of the prescribed field form.
\end{theorem}

\begin{proof}
Let $H\subset\Gamma$ be the preimage of a factor
stabilizer. It has index $t$ and corresponds to an etale
resolution cover $Y_H\to Y$. Its action on the first
factor is a character $\chi:H\to\mathcal G_q$ that
extends over $Y_H$, since $\beta$ extends over $Y$.

We spell out the twisted nonabelian Shapiro map. Put
$X=\Gamma/H$. Choose coset representatives $\gamma_x$,
with $\gamma_1=1$, and identify the copy of $S$ at $x$
with the base copy using the given field transport of
$\beta(\gamma_x)$. In these coordinates the $\Gamma$
action on $A=S^X$ is
\[
 (\gamma a)_x=\chi(h(\gamma,x))(a_{\gamma^{-1}x}),
 \qquad
 h(\gamma,x)=\gamma_x^{-1}\gamma\gamma_{\gamma^{-1}x}\in H.
                                                        \tag{D.SC.11}
\]
The identity
$h(\gamma\delta,x)=h(\gamma,x)h(\delta,\gamma^{-1}x)$
checks the action. Lifts of $\beta$ to $L$, modulo
$A$-conjugation, are $H^1(\Gamma,A)$ for this action.
For a cocycle $z:H\to S$ with
$z(hh')=z(h)\chi(h)(z(h'))$, the formula
\[
                    c_z(\gamma)_x=z(h(\gamma,x))       \tag{D.SC.12}
\]
is a $\Gamma$-cocycle. Conversely evaluate a cocycle at
the base coset and restrict to $H$. If
$a_x=c(\gamma_x)_x$, the cocycle identity gives
\[
 c_z(\gamma)_x=a_x^{-1}c(\gamma)_x
                      \chi(h(\gamma,x))(a_{\gamma^{-1}x}).
                                                        \tag{D.SC.13}
\]
Thus evaluation and (D.SC.12) induce inverse bijections
on the two nonabelian cohomology sets. All choices have
finite index, so continuity is preserved.

The $k$ coordinate lifts in (D.SC.9) are distinct in
$H^1(\Gamma,A)$. Otherwise an $A$-conjugacy between
two would put their image in a proper conjugate
diagonal over $P$, contrary to surjectivity. Under
(D.SC.12), each class is represented by a homomorphism
$h\mapsto(z(h),\chi(h))$ to
$S\rtimes\mathcal G_q$. The fixed embedding into
$\operatorname{GL}_r(\mathbb F_q)\rtimes\mathcal G_q$
is compatible with its field action. An ambient
$\operatorname{GL}_r(\mathbb F_q)$-conjugacy class
contains at most $|\operatorname{GL}_r(\mathbb F_q)|$
classes under $S$: every conjugating matrix determines
at most one homomorphism. No surjectivity to the ambient
linear group is required for this crude bound. Hence
Theorem~\ref{D:etale-resolution-semilinear-count} gives
\[
 k\leq q^{(A_R't+1)r^2}
       \leq |S|^{\kappa(A_R'+1)t}\leq |A|^c.           \tag{D.SC.14}
\]

Transitivity makes $A$ minimal normal in $L$. The
injectivity of the outer action makes its centralizer
trivial, so $A$ is the unique minimal normal subgroup.
It is nonabelian and is not contained in the Frattini
subgroup. A maximal subgroup not containing $A$ is
core-free; thus $L$ is monolithic primitive, as required
by the crown theorem. Moreover $P$ is a quotient of
$\pi_1(Y)\simeq\pi_1(C)$, so $d(P)\leq2g$.

Apply the already checked crown estimate in
Theorem~\ref{D:unramified-permutation-linear-crowns},
namely \citet[Proposition 2.6, equations (2.1)--(2.2),
pp.~12--13]{LucchiniMarionTracey2020D}. For
$s\geq\max\{2,d(P)\}$ it gives
\[
 \psi_L(s)\geq\frac{|A|^{s-1}}{2t|\operatorname{Out}(S)|}
                  \geq |A|^{s-2},
 \qquad k\leq\psi_L(s)\Longrightarrow d(L_k)\leq s.
\]
Take $s=\max\{2,2g,c+2\}$ and use (D.SC.14).
This proves (D.SC.10), retaining the classification
dependence of that finite-group input.

For a concrete unbounded-rank scope, the existing natural
linear/order bound applies to
$S=\operatorname{SL}_{p^a}(\mathbb F_{p^f})$ whenever
this group is nonabelian simple. Its scalar centre is
trivial since $\gcd(p^a,p^f-1)=1$, it is stable under
all field automorphisms, and the previously checked
order estimate gives (D.SC.8) with $\kappa=4$.
The nonsimple small parameters are excluded by the
stated hypothesis. This example is a scope consequence
of the same crown bound, not another movement event.
\end{proof}
